% Appendix

\section*{Appendix A: Formal Proof of Safe Parallel Integration}

\subsection*{A.1 Preliminaries and Notation}

Let $\mathcal{N} = \langle P, T, \text{Pre}, \text{Post}, m_0, k, S, \Phi, \Sigma, \text{Reg}, \Psi, E, A \rangle$ be a Signal Hierarchical Petri Net. For transitions $t_1, t_2 \in T$, we denote:
\begin{itemize}
\item ${}^\bullet t_i = \{p \in P \mid \text{Pre}(p, t_i) > 0\}$ (input places)
\item $t_i^\bullet = \{p \in P \mid \text{Post}(t_i, p) > 0\}$ (output places)
\item $\Sigma(t_i) = \{p \in P \mid p \text{ connected to } t_i \text{ via test/inhibitor arc}\}$ (regulatory places)
\item $\Phi(t_i, M)$ (rate function dependent on marking $M$)
\end{itemize}

\subsection*{A.2 Lemma 1: ODE Superposition for Convergent Coupling}

\begin{lemma}[Convergent Superposition]
For weakly independent transitions $t_1, t_2$ with shared output place $p \in t_1^\bullet \cap t_2^\bullet$ and ${}^\bullet t_1 \cap {}^\bullet t_2 = \emptyset$, parallel integration yields:
\[
M(p, t + \Delta t) = M(p, t) + \int_t^{t+\Delta t} \left[\Phi(t_1, M(\tau)) \cdot \text{Post}(t_1, p) + \Phi(t_2, M(\tau)) \cdot \text{Post}(t_2, p)\right] d\tau
\]
This equals sequential integration order (either $t_1 \to t_2$ or $t_2 \to t_1$).
\end{lemma}

\begin{proof}
By linearity of integration and commutativity of addition:
\begin{align*}
\int_t^{t+\Delta t} \left[\Phi(t_1) + \Phi(t_2)\right] d\tau &= \int_t^{t+\Delta t} \Phi(t_1) d\tau + \int_t^{t+\Delta t} \Phi(t_2) d\tau \\
&= \lim_{n \to \infty} \sum_{i=1}^n \Phi(t_1, \tau_i) \Delta\tau + \lim_{m \to \infty} \sum_{j=1}^m \Phi(t_2, \tau_j) \Delta\tau
\end{align*}

Since ${}^\bullet t_1 \cap {}^\bullet t_2 = \emptyset$, the rate functions $\Phi(t_1, M)$ and $\Phi(t_2, M)$ depend on disjoint substrate sets. The order of evaluation does not affect the result. Sequential execution in either order yields identical marking evolution.
\end{proof}

\subsection*{A.3 Lemma 2: Read-Only Regulatory Coupling}

\begin{lemma}[Regulatory Reads are Commutative]
For weakly independent transitions $t_1, t_2$ with shared regulatory place $p \in \Sigma(t_1) \cap \Sigma(t_2)$ and ${}^\bullet t_1 \cap {}^\bullet t_2 = \emptyset$, simultaneous evaluation of $\text{Reg}(p, M(t))$ yields identical results to sequential evaluation.
\end{lemma}

\begin{proof}
Test arcs and inhibitor arcs read $M(p)$ without modification:
\[
\text{Reg}(p, M(t)) = \begin{cases}
M(p) & \text{(test arc)} \\
\max(0, \theta_p - M(p)) & \text{(inhibitor arc)}
\end{cases}
\]

Both transitions evaluate $\text{Reg}(p, M(t))$ at the same marking $M(t)$. No race condition exists because $M(p)$ is not consumed during regulatory evaluation. Parallel reads commute:
\[
\text{Reg}_1(\text{Reg}_2(M)) = \text{Reg}_2(\text{Reg}_1(M)) = M
\]
\end{proof}

\subsection*{A.4 Lemma 3: Disjoint Substrate Independence}

\begin{lemma}[Disjoint Substrates Guarantee Independence]
For transitions $t_1, t_2$ with ${}^\bullet t_1 \cap {}^\bullet t_2 = \emptyset$, substrate consumption operations commute:
\[
M_{\text{parallel}} = M - \text{Pre}(\cdot, t_1) - \text{Pre}(\cdot, t_2) = M_{\text{seq-12}} = M_{\text{seq-21}}
\]
\end{lemma}

\begin{proof}
Substrate consumption for transition $t_i$ modifies marking:
\[
M'(p) = M(p) - \text{Pre}(p, t_i) \cdot \delta_{p \in {}^\bullet t_i}
\]

Since ${}^\bullet t_1 \cap {}^\bullet t_2 = \emptyset$, the sets of modified places are disjoint. Vector subtraction is commutative for disjoint index sets:
\begin{align*}
M_{\text{parallel}} &= M - \sum_{p \in {}^\bullet t_1} \text{Pre}(p, t_1) \cdot e_p - \sum_{p \in {}^\bullet t_2} \text{Pre}(p, t_2) \cdot e_p \\
&= M - \sum_{p \in {}^\bullet t_2} \text{Pre}(p, t_2) \cdot e_p - \sum_{p \in {}^\bullet t_1} \text{Pre}(p, t_1) \cdot e_p \\
&= M_{\text{seq-21}}
\end{align*}
where $e_p$ denotes the unit vector for place $p$.
\end{proof}

\subsection*{A.5 Theorem: Safe Parallel Integration (Complete)}

\begin{theorem}[Safe Parallel Integration]
For continuous transitions $t_1, t_2$ that are weakly independent (${}^\bullet t_1 \cap {}^\bullet t_2 = \emptyset$), parallel ODE integration over time interval $[t, t+\Delta t]$ preserves correctness: the resulting marking $M(t+\Delta t)$ is identical to sequential integration in any order.
\end{theorem}

\begin{proof}
Consider the marking evolution operator $\Theta_{\Delta t}(M, T)$ for transition set $T$ over interval $\Delta t$. We must show:
\[
\Theta_{\Delta t}(M, \{t_1, t_2\}_{\text{parallel}}) = \Theta_{\Delta t}(M, \{t_1, t_2\}_{\text{seq-12}}) = \Theta_{\Delta t}(M, \{t_1, t_2\}_{\text{seq-21}})
\]

\textbf{Step 1: Decompose marking evolution}

For each place $p \in P$, the marking evolution decomposes into:
\begin{align*}
M(p, t+\Delta t) = M(p, t) &+ \underbrace{\sum_{t_i \in \{t_1, t_2\}} \int_t^{t+\Delta t} \Phi(t_i, M(\tau)) \cdot \text{Post}(t_i, p) d\tau}_{\text{production}} \\
&- \underbrace{\sum_{t_i \in \{t_1, t_2\}} \int_t^{t+\Delta t} \Phi(t_i, M(\tau)) \cdot \text{Pre}(p, t_i) d\tau}_{\text{consumption}}
\end{align*}

\textbf{Step 2: Apply weak independence constraint}

Since ${}^\bullet t_1 \cap {}^\bullet t_2 = \emptyset$, for any place $p$:
\begin{itemize}
\item If $p \in {}^\bullet t_1$, then $p \notin {}^\bullet t_2$, so $\text{Pre}(p, t_2) = 0$
\item If $p \in {}^\bullet t_2$, then $p \notin {}^\bullet t_1$, so $\text{Pre}(p, t_1) = 0$
\end{itemize}

Therefore, consumption terms never overlap. The rate function $\Phi(t_i, M)$ depends only on $M({}^\bullet t_i)$, which is disjoint.

\textbf{Step 3: Analyze coupling modes}

\textit{Case 1: Convergent coupling} ($t_1^\bullet \cap t_2^\bullet \neq \emptyset$). For shared output $p \in t_1^\bullet \cap t_2^\bullet$:
\[
M(p, t+\Delta t) = M(p, t) + \int_t^{t+\Delta t} \left[\Phi(t_1) \cdot \text{Post}(t_1, p) + \Phi(t_2) \cdot \text{Post}(t_2, p)\right] d\tau
\]
By Lemma 1 (ODE superposition), this equals sequential integration.

\textit{Case 2: Regulatory coupling} ($\Sigma(t_1) \cap \Sigma(t_2) \neq \emptyset$). For shared regulatory place $p \in \Sigma(t_1) \cap \Sigma(t_2)$:

Both $\Phi(t_1, M)$ and $\Phi(t_2, M)$ depend on $M(p)$ via $\text{Reg}(p)$, but do not modify $M(p)$. By Lemma 2 (read-only commutation), parallel evaluation yields identical results.

\textit{Case 3: Disjoint} ($t_1^\bullet \cap t_2^\bullet = \emptyset$ and $\Sigma(t_1) \cap \Sigma(t_2) = \emptyset$). No shared places exist. By Lemma 3, operations commute trivially.

\textbf{Step 4: Conclusion}

In all cases (convergent, regulatory, disjoint), the weak independence constraint ${}^\bullet t_1 \cap {}^\bullet t_2 = \emptyset$ ensures that parallel integration yields results identical to sequential integration in any order. Therefore, parallel execution preserves correctness.
\end{proof}

\subsection*{A.6 Corollary: Stochastic Tau-Leaping Parallelism}

\begin{corollary}[Parallel Tau-Leaping]
For stochastic transitions $t_1, t_2$ that are weakly independent, concurrent Poisson sampling within time leap $[t, t+\tau]$ yields statistically equivalent results to sequential sampling.
\end{corollary}

\begin{proof}
Tau-leaping approximates the number of firings $k_i \sim \text{Poisson}(\Phi(t_i) \cdot \tau)$ for each transition. For weakly independent $t_1, t_2$:

\textbf{Independence of Poisson processes:} Since ${}^\bullet t_1 \cap {}^\bullet t_2 = \emptyset$, the propensity functions depend on disjoint substrate sets. The Poisson processes are independent:
\[
P(k_1, k_2) = P(k_1) \cdot P(k_2) = \frac{(\lambda_1 \tau)^{k_1} e^{-\lambda_1 \tau}}{k_1!} \cdot \frac{(\lambda_2 \tau)^{k_2} e^{-\lambda_2 \tau}}{k_2!}
\]

\textbf{Commutativity of marking updates:} The marking update after firing $k_1$ times for $t_1$ and $k_2$ times for $t_2$ is:
\[
M' = M - k_1 \cdot \text{Pre}(\cdot, t_1) - k_2 \cdot \text{Pre}(\cdot, t_2) + k_1 \cdot \text{Post}(t_1, \cdot) + k_2 \cdot \text{Post}(t_2, \cdot)
\]

Since pre-sets are disjoint, the order of updates does not affect $M'$. Parallel sampling yields statistically equivalent results.
\end{proof}

\subsection*{A.7 Complexity Analysis}

\begin{proposition}[Classification Complexity]
Given a Petri net with $|T|$ transitions and $|F|$ arcs, pairwise weak independence classification requires $O(|T|^2 \cdot |F|)$ time.
\end{proposition}

\begin{proof}
For each pair $(t_1, t_2)$, we must check ${}^\bullet t_1 \cap {}^\bullet t_2 = \emptyset$:
\begin{itemize}
\item Computing ${}^\bullet t_i$ requires scanning incoming arcs: $O(|F|/|T|)$ per transition
\item Checking intersection of two pre-sets: $O(|P|)$ worst case
\item Total for one pair: $O(|F|/|T| + |P|) = O(|F|)$ (since $|P| \leq |F|$)
\item Total for all $\binom{|T|}{2} = O(|T|^2)$ pairs: $O(|T|^2 \cdot |F|)$
\end{itemize}

This is efficient for typical biological networks where $|T| \ll |F|$.
\end{proof}
